\documentclass{article}
\usepackage[UTF8]{ctex}
\usepackage{amsmath,amssymb,amsfonts}
\usepackage{enumitem}
\usepackage{geometry}
\geometry{a4paper, margin=1in}
% === 错题本排版优化（不参与题目标识） ===
\setlength{\parindent}{0pt}
\setlength{\parskip}{0.35em plus 0.10em minus 0.08em}
\setlength{\emergencystretch}{3em}
\setlist{itemsep=0.25em, topsep=0.40em, parsep=0pt}
\allowdisplaybreaks[2]
% === 错题本排版优化结束 ===
\begin{document}
\section*{中国科学技术大学2025年高等代数真题}
\section*{98．中国科学技术大学2025年线性代数与解析几何试卷}
一．填空题。\\

\subsection{题目1}
1．在空间直角坐标系中，设 \(A=(1,1,1), B=(1,2,3), C=(2,3,4), D=(3,3,4)\) ，则平面 \(A B C\) 的方程为 \(\_\_\_\_\)，\(\triangle A B C\) 的面积为 \(\_\_\_\_\)，四面体 \(A B C D\) 的体积为 \(\_\_\_\_\)．

\subsection{题目2}
2．点 \(A(1,2,3)\) 到 \(l: \frac{x-5}{-4}=\frac{y-4}{3}=\frac{z-3}{-2}\) 的距离为 \(\_\_\_\_\)．

\subsection{题目3}
3．设 \(A=\left(\begin{array}{cccc}1 & 1 & 1 & 1 \\ -1 & 1 & -1 & 1 \\ -1 & -1 & 1 & 1 \\ -1 & -1 & -1 & 1\end{array}\right)\) ，则 \(A^{2025}=\) \(\_\_\_\_\)，\(A\) 的特征多形式 \(\varphi(x)=\) \(\_\_\_\_\)．\\
\(\_\_\_\_\)

\subsection{题目4}
4．设实方阵 \(A=J_{2025}(0)\) 是 Jordan 块，则 \(A^{10}\) 的 Jordan 标准形为 \(\_\_\_\_\)．

\subsection{题目5}
5． \(\mathbb{R}[x]\) 上的多项式矩阵

\[
\left(\begin{array}{ccc}
x-1 & x^{2}-1 & x^{3}-1 \\
x^{2}-1 & x^{3}-1 & x^{4}-1 \\
x^{3}-1 & x^{4}-1 & x^{5}-1
\end{array}\right)
\]

的 Smith 标准型为 \(\_\_\_\_\)．

\subsection{题目6}
6．设 \(p, q\) 分别是实二次型

\[
Q\left(x_{1}, x_{2}, x_{3}, x_{4}, x_{5}\right)=x_{1} x_{2}+x_{2} x_{3}+x_{3} x_{4}+x_{4} x_{5}
\]

的正负惯性指数，则 \((p, q)=\) \(\_\_\_\_\)．

\subsection{题目7}
7．设实线性空间 \(V\) 中的向量组 \(\alpha_{1}, \alpha_{2}, \alpha_{3}, \alpha_{4}, \alpha_{5}\) 线性无关，则向量组

\[
\alpha_{1}+\alpha_{2}, \alpha_{2}+\alpha_{3}, \alpha_{3}+\alpha_{4}, \alpha_{4}+\alpha_{5}, \alpha_{5}+\alpha_{1}
\]

的秩为 \(\_\_\_\_\)．

二．简答题．\\

\subsection{题目8}
1．\(A, B\) 为 \(n\) 阶复方阵，证明： \(\operatorname{det}\left(\begin{array}{cc}A & -B \\ B & A\end{array}\right)=\operatorname{det}(A+\mathrm{i} B) \operatorname{det}(A-\mathrm{i} B)\) ，其中 i 为虚数单位。

\subsection{题目9}
2．\(A, B\) 为 \(n\) 阶实方阵，证明：线性方程组 \(A X=0\) 和 \(B X=0\) 有相同的解集的充要条件为

\[
\operatorname{rank}\binom{A}{B}=\operatorname{rank}(A)=\operatorname{rank}(B) .
\]

\subsection{题目10}
3．\(A\) 为实方阵，证明：若存在实方阵 \(P\) ，使得 \(A^{\mathrm{T}}=A P\) ，则存在正交方阵 \(Q\) ，使得 \(Q^{-1} A Q\) 形如

\[
\left(\begin{array}{ll}
B & O \\
O & O
\end{array}\right)
\]

\(B\) 为可逆方阵．

\subsection{题目11}
4．设实线性空间 \(V\) 上的线性变换 \(\mathcal{A}\) 满足 \(\mathcal{A}^{3}=\mathcal{A}, \mathcal{E}\) 是 \(V\) 上的恒等变换．证明：

\[
V=\operatorname{Im}\left(\mathcal{A}^{2}+\mathcal{A}\right) \oplus \operatorname{Im}\left(\mathcal{A}^{2}-\mathcal{A}\right) \oplus \operatorname{Im}\left(\mathcal{A}^{2}-\mathcal{E}\right) .
\]

\subsection{题目12}
5．设欧氏空间 \(V=\left\{a x^{2}+b x+c \mid a, b, c \in \mathbb{R}\right\}\) 上的内积 \((f, g)=\int_{0}^{1} f(x) g(x) \mathrm{d} x\) 。求 \(V\) 上的线性变换 \(\mathcal{A}: f(x) \rightarrow f(1-x)\) 的伴随变换 \(\mathcal{A}^{*}\) 在基 \(1, x, x^{2}\) 下的矩阵。

\subsection{题目13}
6．求所有的实数 \(a\) ，使得方程 \(x^{2}+y^{2}+z^{2}+x y+x z+a y z=1\) 表示一个圆柱面，并求此圆柱面半径．
\end{document}
